% Slides — Capítulo 15: Perigos de Tempestades
\documentclass[aspectratio=169,11pt]{beamer}
\usepackage{../comum/temameteo}
\graphicspath{{../figuras/cap15/}}

\title[Cap. 15 --- Perigos]{Capítulo 15 --- Perigos de Tempestades}
\subtitle{Meteorologia Prática}
\author{}
\date{}

\begin{document}

\begin{frame}
  \titlepage
  \vfill
  \atribuicao
\end{frame}

\begin{frame}{Roteiro}
  \tableofcontents
\end{frame}

% ---------------------------------------------------------------
\section{Granizo}

\begin{frame}{A balança updraft × queda}
  \begin{columns}
    \column{0.4\textwidth}
    \eqdestaque{$w_T = \sqrt{\dfrac{4g\rho_{gelo}D}{3C_d\rho_{ar}}}
      \propto \sqrt{D}$}
    \vspace{0.4em}
    \begin{itemize}
      \item Pedra de 5\,cm exige updraft $\gtrsim$32\,m/s: granizo
            grande \alert{denuncia} supercélula (Cap.~14).
      \item As camadas de cebola: voos repetidos pela zona de
            crescimento (riming, Cap.~7).
    \end{itemize}
    \column{0.3\textwidth}
    \figlivroslide{fig_15_4.jpg}{0.42\textheight}{Fig.~15.4 ---
      velocidade terminal vs.\ diâmetro}
    \column{0.3\textwidth}
    \figlivroslide{fig_15_6.jpg}{0.42\textheight}{Fig.~15.6 --- o
      embrião de graupel e as camadas}
  \end{columns}
\end{frame}

\begin{frame}{Onde a pedra cresce}
  \begin{columns}
    \column{0.4\textwidth}
    \begin{itemize}
      \item Zona ótima: $-10$ a $-30\,°$C, com muita água
            super-resfriada (Cap.~7).
      \item BWER no radar: o updraft tão forte que nem a chuva cai
            dentro (Cap.~8).
      \item Derretimento na queda: nível de congelamento alto poupa o
            chão (N1) — por isso o equador tem pouco granizo.
    \end{itemize}
    \column{0.6\textwidth}
    \figlivroslide{fig_15_10.jpg}{0.54\textheight}{Fig.~15.10 --- a
      trajetória de crescimento e o BWER}
  \end{columns}
\end{frame}

% ---------------------------------------------------------------
\section{Downbursts}

\begin{frame}{O CAPE de cabeça para baixo}
  \begin{columns}
    \column{0.34\textwidth}
    \eqdestaque{$w_{down} = \sqrt{2\,\mathrm{DCAPE}}$}
    \vspace{0.4em}
    \begin{itemize}
      \item Evaporação esfria a parcela, que desce pelo seu
            $\theta_w$ (Cap.~4).
      \item DCAPE 500--1500\,J/kg $\to$ rajadas de 30--55\,m/s.
    \end{itemize}
    \column{0.33\textwidth}
    \figlivroslide{fig_15_7.jpg}{0.44\textheight}{Fig.~15.7 --- DCAPE
      no Skew-$T$: a área do downdraft}
    \column{0.33\textwidth}
    \figlivroslide{fig_15_14.jpg}{0.44\textheight}{Fig.~15.14 ---
      $w_{down}$ vs.\ DCAPE}
  \end{columns}
\end{frame}

\begin{frame}{No chão: microburst e frente de rajada}
  \begin{columns}
    \column{0.4\textwidth}
    \begin{itemize}
      \item O jato desce, espalha e enrola: vórtice de anel na frente
            de rajada (Cap.~14).
      \item Downburst \alert{seco}: cai de virga, quase sem eco — a
            armadilha da aviação (N2).
      \item Haboob: a parede de poeira da rajada em solo seco.
    \end{itemize}
    \column{0.6\textwidth}
    \figlivroslide{fig_15_11.jpg}{0.52\textheight}{Fig.~15.11 ---
      downburst, ventos de outflow e a frente de rajada}
  \end{columns}
\end{frame}

% ---------------------------------------------------------------
\section{Relâmpagos}

\begin{frame}{O gerador: colisões graupel--cristal}
  \begin{columns}
    \column{0.34\textwidth}
    \begin{itemize}
      \item Graupel (pesado) fica $-$, cristais (leves) sobem $+$:
            o updraft separa cargas.
      \item Exige gelo $+$ água super-resfriada $+$ updraft: raio e
            granizo são irmãos (Cap.~7).
      \item Dente de serra do capacitor: Caso~3 do notebook.
    \end{itemize}
    \column{0.33\textwidth}
    \figlivroslide{fig_15_22b.jpg}{0.42\textheight}{Fig.~15.22 --- a
      transferência de carga nas colisões}
    \column{0.33\textwidth}
    \figlivroslide{fig_15_23.jpg}{0.42\textheight}{Fig.~15.23 --- líder
      escalonado e descarga de retorno}
  \end{columns}
\end{frame}

\begin{frame}{Trovão e segurança}
  \begin{columns}
    \column{0.4\textwidth}
    \begin{itemize}
      \item Canal a $\sim$30\,000\,K: expansão supersônica $=$
            trovão.
      \item \alert{3\,s por km} (flash-to-bang, T3); regra 30--30
            para abrigo.
      \item O som refrata para cima (Cap.~1: $T$ cai com $z$): além
            de $\sim$15\,km o raio é ``mudo''.
    \end{itemize}
    \column{0.6\textwidth}
    \figlivroslide{fig_15_31.jpg}{0.5\textheight}{Fig.~15.31 ---
      refração do trovão e o alcance audível}
  \end{columns}
\end{frame}

% ---------------------------------------------------------------
\section{Tornados}

\begin{frame}{O zoológico e o ciclo de vida}
  \begin{columns}
    \column{0.3\textwidth}
    \begin{itemize}
      \item Formas: cone, cunha, ampulheta, corda — tamanho não é
            força.
      \item Família: supercelular, landspout, waterspout, gustnado
            (e o primo dust devil).
      \item Estágios: organização $\to$ maduro $\to$ encolhimento
            $\to$ corda.
    \end{itemize}
    \column{0.35\textwidth}
    \figlivroslide{fig_15_32b.jpg}{0.42\textheight}{Fig.~15.32 ---
      formas de funil}
    \column{0.35\textwidth}
    \figlivroslide{fig_15_38.jpg}{0.42\textheight}{Fig.~15.38 --- o
      ciclo de vida}
  \end{columns}
\end{frame}

\begin{frame}{Rankine: o vórtice e a queda de pressão}
  \begin{columns}
    \column{0.4\textwidth}
    \eqdestaque{$\dfrac{dP}{dr} = \dfrac{\rho v^2}{r}
      \Rightarrow \Delta P = \rho v_{max}^2$}
    \vspace{0.4em}
    \begin{itemize}
      \item Ciclostrófico puro (Cap.~10): sem Coriolis, gira para
            qualquer lado.
      \item EF5 (110\,m/s): $\sim$120\,hPa em 100\,m.
      \item Carga $q = \tfrac12\rho v^2$: dobrar $v$ quadruplica a
            força (T5).
    \end{itemize}
    \column{0.6\textwidth}
    \figlivroslide{fig_15_33.jpg}{0.52\textheight}{Fig.~15.33 --- o
      perfil de Rankine e a depressão central}
  \end{columns}
\end{frame}

\begin{frame}{De onde vem o giro}
  \begin{columns}
    \column{0.34\textwidth}
    \begin{itemize}
      \item Vorticidade horizontal do cisalhamento é
            \alert{inclinada} pelo updraft e \alert{esticada}
            (Cap.~13!) até virar tornado.
      \item Mesociclone $\to$ RFD $\to$ toque no solo.
      \item SRH do hodógrafo prevê o risco (Cap.~14, N3).
    \end{itemize}
    \column{0.33\textwidth}
    \figlivroslide{fig_15_41a.jpg}{0.42\textheight}{Fig.~15.41 ---
      inclinar e esticar o tubo de vorticidade}
    \column{0.33\textwidth}
    \figlivroslide{fig_15_44a.jpg}{0.42\textheight}{Fig.~15.44 --- do
      mesociclone ao vórtice de tornado}
  \end{columns}
\end{frame}

\begin{frame}{Translação e vórtices de sucção}
  \begin{columns}
    \column{0.34\textwidth}
    \begin{itemize}
      \item Giro $+$ translação: no HS o lado \alert{esquerdo} da
            trajetória é o pior (N4).
      \item Tornados multivórtice: sub-vórtices de sucção varrem
            faixas estreitas de dano extremo.
      \item Escala EF: dano $\to$ vento estimado.
    \end{itemize}
    \column{0.33\textwidth}
    \figlivroslide{fig_15_34.jpg}{0.42\textheight}{Fig.~15.34 ---
      vento total $=$ rotação $+$ translação}
    \column{0.33\textwidth}
    \figlivroslide{fig_15_52.jpg}{0.42\textheight}{Fig.~15.52 ---
      vórtices de sucção múltiplos}
  \end{columns}
\end{frame}

% ---------------------------------------------------------------
\section{Síntese e laboratório}

\begin{frame}{O mapa do capítulo}
  \eqdestaque{$w_T \propto \sqrt{D}$ \quad
    $w_{down} = \sqrt{2\mathrm{DCAPE}}$ \quad
    $\dot Q = I_c - Q/RC$ \quad $\Delta P = \rho v_{max}^2$}
  \vspace{0.6em}
  Quatro perigos, quatro pedaços de física clássica no extremo — todos
  filhos da mesma supercélula do Cap.~14, todos visíveis no radar do
  Cap.~8.
\end{frame}

\begin{frame}{Laboratório numérico --- notebook do Cap.~15}
  \texttt{notebooks/cap15\_perigos.ipynb}
  \begin{enumerate}
    \item Granizo: $w_T(D)$ com os marcos ervilha$\to$laranja.
    \item Downburst: a EDO da queda com déficit evaporativo.
    \item O capacitor da tempestade (dente de serra de raios).
    \item Rankine: $v(r)$ e $\Delta P$ por categoria EF.
  \end{enumerate}
  \vspace{0.4em}
  \textbf{Exercícios}: T1--T5 e N1--N4 nas notas; células reservadas no
  notebook.
  \vfill
  \atribuicao
\end{frame}

\end{document}
